\section{A thermal slab and joint trajectory estimates}\label{sec:slab}

The spatial estimates control the equilibrium positions.  We now need
the corresponding control of trajectories on a short time interval,
simultaneously near any prescribed set of cells.  The first step is an
exact path representation of the Gibbs state that also retains the
off-diagonal information needed for condensation.  The second step
transfers rare events for a free path to joint estimates for the
interacting sample.  For this transfer we adapt the auxiliary-color
and trace-sandwich argument of \citet[Section ``A joint bound for
atypical trajectories,'' in particular the lemma ``Trace bound for a
Lipschitz sandwich'']{B}.  The hard-core domains and the trajectory
tests used here require the details given below.

\subsection{The path law and its purification}

Fix $0<t\le1$, set $\delta=2t$, and require $\beta>\delta$.
The time $t$ will be fixed before the lower threshold on $D$.
Let $\nu_N$ be the measure of $N$ independent torus Brownian motions
with generator $\Delta$ on $[0,\delta]$, with Lebesgue initial
measure.  Write
\[
 X(s)=(\omega_1(s),\ldots,\omega_N(s)),\qquad
 Y=X(0),\qquad Z=X(\delta).
\]
In addition to these labelled paths, introduce a configuration
$W'\in\Lambda_K^N$.  With the symmetric heat-kernel convention of
Section~\ref{sec:setup}, define
\begin{equation}\label{eq:thermal-slab-law}
 \dd\cR(\omega,W')=
 \frac{B_{\beta/2-t}^{\rm sym}(W',Y)
       B_{\beta/2-t}^{\rm sym}(Z,W')}{Z_N}
 \ind_{\{d(\omega_i(s),\omega_j(s))>r\ \forall s,\ i<j\}}
 \dd\nu_N(\omega)\,\dd W'.
\end{equation}
All path spaces carry their Borel sigma fields.  Continuous torus
paths have continuous Euclidean lifts, unique after specifying the
initial lift; lifted increments and diameters are independent of that
choice.

\begin{proposition}[Thermal purification]\label{prop:purification}
Equation~\eqref{eq:thermal-slab-law} defines a probability measure.
At every fixed time $s\in[0,\delta]$, its position marginal is
\[
 \mu(\dd X)=Z_N^{-1}B_\beta^{\rm sym}(X,X)\,\dd X.
\]
The joint law of $X(t)$ and $W'$ has density $\Phi(X,W')^2$, where
\begin{equation}\label{eq:thermal-Phi}
 \Phi(X,W')=Z_N^{-1/2}B_{\beta/2}^{\rm sym}(X,W').
\end{equation}
This function is nonnegative, has squared integral one, and is
symmetric separately in $X$ and $W'$.  The constant-orbital
occupation fraction of the scaled Gibbs state is
\begin{equation}\label{eq:occupation-purification}
 B_{\rm occ}
 =\frac1V\int\left|\int_{\Lambda_K}
       \Phi(x,X',W')\,\dd x\right|^2\dd X'\,\dd W',
\end{equation}
where $X'\in\Lambda_K^{N-1}$ is an ordered configuration.
\end{proposition}

\begin{proof}
The avoidance indicator gives the labelled Dirichlet heat kernel by
the open-domain Feynman--Kac representation established in
Section~\ref{sec:setup}, with the same strict-avoidance and zero-extension
conventions.

Let $P_{\rm sym}$ be the orthogonal symmetrizer in labelled
configuration space.  It commutes with the labelled semigroup and
$P_{\rm sym}^2=P_{\rm sym}$.  Integrating $W'$ in
Equation~\eqref{eq:thermal-slab-law} therefore leaves the outer
kernel $B_{\beta-\delta}^{\rm sym}(Z,Y)$.  Composing with the
slab kernel and taking the trace gives $Z_N$, so the total mass
is one.  Splitting the labelled slab kernel at time $s$ gives
the diagonal kernel $B_\beta^{\rm sym}(X,X)$ at that time.
These identities hold against arbitrary bounded position tests;
alternatively all diagonal traces can be evaluated after a further
positive-time split.

If instead we retain $W'$ and split at $t$, composing the first
outer kernel with the first killed half-slab gives
$B_{\beta/2}^{\rm sym}(W',X)$.  The other half gives its transpose.
Their product, divided by $Z_N$, is $\Phi(X,W')^2$.

Finally $\Phi$ is the zero-extended kernel of $\Gamma_N^{1/2}$.
Thus the kernel of the trace-one Gibbs state satisfies
\[
 \Gamma_N(X,\bar X)=\int
       \Phi(X,W')\Phi(\bar X,W')\,\dd W'.
\]
The reduced density matrix is $N$ times the partial trace over
$X'$.  Its expectation in the orbital $V^{-1/2}$, divided by $N$,
is precisely Equation~\eqref{eq:occupation-purification}.
Tonelli's theorem applies since the kernels are nonnegative;
the resulting integral is finite since a one-particle occupation
fraction is at most one.
\end{proof}

The extra variable $W'$ is what permits a mixed Gibbs state to be
treated by a nonnegative squared density.  It will be held fixed
in the later path changes.  No ground-state eigenfunction enters
Equation~\eqref{eq:thermal-slab-law}.

\subsection{Regular trajectories and certificates}

Recall $s_0=bD^{-1/3}$ and the fixed parameters
$\ell,\theta,b$ from the spatial estimates, with
$3\theta<\ell/2$.  Put
\begin{equation}\label{eq:regularity-scales}
 h_0=D^{-1},\qquad \gamma=\frac1{12},\qquad
 m_t=\lceil t/r^2\rceil,\qquad \tau=t/m_t.
\end{equation}
Take $D$ large enough that $h_0<t$ and
$r^2/2\le\tau\le r^2$.  Divide each half of the slab into
$m_t$ equal intervals, denoted together by $I_1,\ldots,I_{2m_t}$.
For a lifted path define
\[
 A_j(\omega)=r^{-1}\diam\bigl(\widetilde\omega(I_j)\bigr).
\]
The following definition is the one used in
\citet[Section ``A stationary sample of hard-core paths,''
definition of a regular path]{HS}.

\begin{definition}[Regular path]\label{def:regular}
Fix a sufficiently large numerical constant $C_{\rm reg}$.
A path on $[0,2t]$ is \emph{regular} if
\begin{equation}\label{eq:regularity-tests}
 \max_j A_j\le D^\gamma,
 \qquad
 \frac1{2m_t}\sum_{j=1}^{2m_t}(1+A_j)^3\le C_{\rm reg},
\end{equation}
and its lifted diameter in each endpoint window
$[0,h_0]$ and $[2t-h_0,2t]$ is at most $s_0/8$.
\end{definition}

We will use both the free-motion and the bridge forms of the next
estimate.  The latter supplies the individual regularity tests
for the replacement paths in Section~\ref{sec:bridges}.

\begin{lemma}[Free cutoff bounds]\label{lem:free-cutoffs}
Fix $t>0$.  For every $q>0$, a free Brownian path on $[0,2t]$
fails regularity with probability at most $C_{q,t}D^{-q}$,
uniformly in its starting point, $K$, and the allowed
$\alpha$-window.  The same assertion holds for a Euclidean
bridge of duration $2t$ with bounded endpoint displacement, and
for two independent Euclidean bridge legs of duration $t$
joined at a prescribed midpoint, when their three prescribed
points lie in a fixed bounded region.  Constants and the lower
threshold on $D$ may depend on that region and $t$, but
$C_{\rm reg}$ is a single numerical constant for all these laws.
If the three points lie in the radius-$3$ ball about a site $v$,
the probability that the two-leg path leaves the radius-$10$
ball about $v$ is at most $Ce^{-c/t}$ for $0<t\le1$.
\end{lemma}

\begin{proof}
We reproduce the argument of \citet[Lemma ``Free cutoff
bounds'']{HS}.  The variables $A_j$ of Brownian motion depend
on disjoint increments.  Brownian scaling and
$\tau/r^2\in[1/2,1]$ give
$\PP(A_j>a)\le Ce^{-ca^2}$ for $a\ge1$.
Hence the maximum test fails with probability at most
$C m_t e^{-cD^{2\gamma}}$.

Set $U_j=(1+A_j)^3$.  Every fixed moment of $U_j$ is bounded
uniformly.  In the expansion of the $2k$th moment of
$\sum_j(U_j-\E U_j)$, a nonzero term uses every index at
least twice, by independence and centering.  At most $k$
distinct indices can occur, so
\[
 \E\left|\frac1{2m_t}\sum_j(U_j-\E U_j)\right|^{2k}
 \le C_k m_t^{-k}.
\]
Choose $C_{\rm reg}$ larger than a fixed multiple of
$1+\sup_j\E U_j$.  Since $m_t$ is comparable to $D^2$
for fixed $t$ and the fixed $\alpha$-window, Markov's
inequality gives every prescribed inverse power of $D$ for
the average test.  The same argument applies to the average
on either half separately, with the same threshold.
For each endpoint window the maximal Gaussian estimate gives
\[
 \PP\{\diam\widetilde\omega([0,h_0])>s_0/8\}
 \le C\exp(-c s_0^2/h_0)
 =C\exp(-c b^2D^{1/3}),
\]
and reversal handles the other window.

For a bridge leg, write its centered part as
$B(u)-(u/t)B(t)$ and add the line joining its prescribed
endpoints.  On a mesh interval, the additional increment in
units of $r$ is bounded by
$C(r/t)(|B(t)|+L)$, where $L$ bounds endpoint displacement.
On the event $|B(t)|\le D^{1/12}$ this tends to zero,
while the complementary probability is at most
$C_t e^{-c_tD^{1/6}}$.  The maximum test follows by applying
the Brownian bound with threshold $D^\gamma/2$.
The inequality $(x+y)^3\le4x^3+4y^3$, with a sufficiently
large fixed choice of $C_{\rm reg}$, transfers the average
test as well.  The additional displacement on an endpoint
window is at most $C h_0(|B(t)|+L)/t=o(s_0)$ on the same
event; use the Brownian window estimate with half the allowed
threshold.  Two independent legs require a union bound, and
a duration-$2t$ bridge has the identical representation with
$2t$ in place of $t$.

For the final assertion, each interpolating line remains in
the radius-$3$ ball.  Exiting the radius-$10$ ball requires
a centered fluctuation of size at least seven on one leg.
The Brownian representation and its maximal Gaussian bound
give $Ce^{-c/t}$.
\end{proof}

Write $C_v(X)$ for the close count from
Proposition~\ref{prop:close-count}, namely
\[
 C_v(X)=\#\{i:d(X_i,v)\le5,\ d_i(X)\le s_0\}.
\]
A \emph{certificate at $v$} consists of the following conditions
on the full path sample:
\begin{enumerate}
 \item $n_v(Y)\ge\ell D$ and $C_v(Y),C_v(Z)<\theta D$;
 \item at most $H_1D$ paths visit the radius-$12$ ball about
 $v$, every visiting path is regular, and fewer than
 $\theta D$ visiting paths have lifted diameter greater than one.
\end{enumerate}
These are precisely the spatial and trajectory conditions needed
for the deterministic bridge estimates.  In particular, the
certificate is an event of the full sample; it is not information
on which we will condition the paths to be changed.

\begin{proposition}[Joint certificate supply]\label{prop:certificates}
For each prescribed $p_1>0$, one can choose $H_1$ large,
then $t>0$ sufficiently small, and then $D_0$ sufficiently
large, such that, for every set $A$ of distinct cell centers,
\begin{equation}\label{eq:joint-certificates}
 \cR\{\text{no site of }A\text{ has a certificate}\}
 \le p_1^{|A|}
\end{equation}
whenever $D\ge D_0$, $0<T\le D^{-1}$, and $K$ is
sufficiently large.  The choices are uniform over the
$\alpha$-window.  After $H_1$ is fixed, $t$ may be decreased
further before increasing $D_0$.
\end{proposition}

The endpoint parts already follow from
Propositions~\ref{prop:low-count} and \ref{prop:close-count}.
The issue is to control many path failures jointly under
$\cR$.  We establish that estimate next, and prove
Proposition~\ref{prop:certificates} at the end of the section.

\subsection{A localized kernel for one marked path}

For a nonempty set $A$ of sites let $J_A(\omega)$ be the
number of sites in $A$ whose radius-$12$ balls the path visits.
Consider the following three single-path events:
\[
 E_{\rm all}=\{\text{all paths}\},\qquad
 E_{\rm diam}=\{\diam\widetilde\omega>1\},\qquad
 E_{\rm irr}=\{\omega\text{ is not regular}\}.
\]
For $E$ equal to one of these events and $\lambda>0$, we will bound
\begin{equation}\label{eq:marked-moment-definition}
 \begin{split}
 \mathfrak M
 &=\E_{\cR}\exp\left(
   \lambda\sum_{i=1}^N J_A(\omega_i)\ind_E(\omega_i)\right)\\
 &=\E_{\cR}\prod_{i=1}^N
   \left[1+\bigl(e^{\lambda J_A(\omega_i)}-1\bigr)
                                      \ind_E(\omega_i)\right].
 \end{split}
\end{equation}
Expanding this product selects subsets of marked labels, each carrying
weight $(e^{\lambda J_A(\omega)}-1)\ind_E(\omega)$.
We first bound the free-path kernel for one such marked label.

We partition paths by a nonnegative integer $n$.  On the first
half, mark the first exit from the unit ball about the starting
point, then the first exit from the unit ball about that exit
point, and continue up to time $t$.  On the second half make
the same construction backwards from the terminal point.
Denote the total number of completed exits on the two halves by
$n(\omega)$, and write $n$ for its value on a fixed bin.
Continuity and uniform continuity on the compact time interval
make this number finite for each path.  Between consecutive
marks the lift lies in a unit ball.  Thus on bin $n$,
\begin{equation}\label{eq:bin-geometry}
 J_A\le C(n+1),\qquad
 |\widetilde\omega(2t)-\widetilde\omega(0)|\le n+2.
\end{equation}
If $J_A>0$, an endpoint is at lifted distance at most
$n+3$ from any visiting point: the first half's
endpoint-to-midpoint displacement is at most $n_1+1$,
and the diameter of the other half is at most $n_2+2$.
Thus both endpoints have torus distance at most $n+15$
from $A$.  By Equation~\eqref{eq:bin-geometry}, the unit
tube about the straight lifted endpoint segment projects
into $\{x:d(x,A)<2n+18\}$, hence into
$\{x:d(x,A)<32(n+1)\}$.

Set $O_n=\{x:d(x,A)<32(n+1)+2\}$.  Choose a Lipschitz
cutoff $0\le p_n\le1$ equal to one when
$d(x,A)\le32(n+1)$ and zero when
$d(x,A)\ge32(n+1)+1$, with $|\nabla p_n|\le C$.
Its support is compact in $O_n$.  These definitions
allow $O_n$ to be the whole torus.  They satisfy
\begin{equation}\label{eq:bin-volume}
 |O_n|\le C(n+1)^3|A|,
 \qquad O_n\text{ is covered by at most }C(n+1)^3|A|
 \text{ unit cells}.
\end{equation}
Here $|O_n|$ is volume and $|A|$ is cardinality.

Let $\mathcal K_n(y,z)$ be the unnormalized free torus
path kernel of duration $\delta$ with the marking weight
and the bin-$n$ restriction.  Thus endpoint integration
against $\mathcal K_n$ is the free-path integral against
that weight, without any hard-core restriction.

\begin{lemma}[Marked-path kernel]\label{lem:marked-kernel}
For a fixed numerical $c_1>1$, which can be taken to be $3$,
\begin{equation}\label{eq:marked-kernel}
 \mathcal K_n(y,z)\le l_n p_n(y)p_n(z)
       [e^{c_1\delta\Delta_{O_n}}](y,z),
 \qquad
 l_n=C\lambda(n+1)e^{C\lambda(n+1)}
       e^{-cn^2/t}\varepsilon^{1/3},
\end{equation}
where $\Delta_{O_n}$ has Dirichlet boundary conditions, or
periodic conditions if $O_n=\Lambda_K$.  The three choices
of $E$ admit respectively
\begin{equation}\label{eq:path-epsilon}
 \varepsilon=1,\qquad
 \varepsilon=Ce^{-c/t},\qquad
 \varepsilon=C_tD^{-12}.
\end{equation}
The constants for the first two choices are uniform for
$0<t\le1$; the irregularity constant may depend on the
fixed $t$ and the previously chosen spatial parameters.
\end{lemma}

\begin{proof}
Fix a lift of $y$ and one lift of $z$.  Omit the marking
factor temporarily but retain $J_A>0$, $E$, and bin $n$.
The resulting Euclidean kernel is bounded by each of
\begin{equation}\label{eq:three-kernel-bounds}
 g_\delta(z-y),\qquad
 Ct^{-3/2}e^{-cn^2/t},\qquad
 Ct^{-3/2}\varepsilon,
 \quad
 g_s(u)=(4\pi s)^{-3/2}e^{-|u|^2/(4s)}.
\end{equation}
The first bound follows by discarding all tests.  To prove
the second, at least one half has at least $n/2$ completed
exits when $n\ge1$.  Bound the opposite half's heat kernel
by $Ct^{-3/2}$ and integrate the indicated half as ordinary
Brownian motion; use reversal for the second half.

For completeness, let $\sigma$ be a Brownian exit time
from the unit ball.  The maximal Gaussian bound gives
$\PP(\sigma\le u)\le Ce^{-c/u}$ for small $u$.
Splitting the expectation at $u=s^{-1/2}$ therefore gives
$\E e^{-s\sigma}\le e^{-c'\sqrt s}$ for sufficiently
large $s$.  Successive exit durations are independent
copies of $\sigma$, by the strong Markov property.
If $Q_t$ is their number completed by $t$, then
\[
 \PP(Q_t\ge k)\le e^{st}(\E e^{-s\sigma})^k.
\]
Taking $\sqrt s$ a sufficiently small fixed multiple of
$k/t$ proves $\PP(Q_t\ge k)\le Ce^{-ck^2/t}$.
When $k/t$ is bounded the same statement follows by
increasing $C$; $n=0$ is covered directly by the heat
kernel supremum.  This proves the second bound in
Equation~\eqref{eq:three-kernel-bounds}.

For the third bound, diameter greater than one forces
one half to have diameter greater than $1/2$, since the
two halves meet.  Its Brownian probability is at most
$Ce^{-c/t}$.  Irregularity forces a failure on at least
one half: an endpoint window, the mesh maximum, or the
half-average with threshold $C_{\rm reg}$.  The proof of
Lemma~\ref{lem:free-cutoffs} bounds each such half failure
by $C_tD^{-12}$.  Once again the opposite half is bounded
by its heat kernel supremum.  This also proves the third
bound for $E_{\rm all}$ with $\varepsilon=1$.

Take the geometric mean of the three bounds.  The surviving
Gaussian is a constant multiple of $g_{3\delta}(z-y)$,
and the other factors become
$e^{-cn^2/t}\varepsilon^{1/3}$ after changing $c,C$.
Equation~\eqref{eq:bin-geometry} bounds the omitted mark by
\[
 e^{\lambda J_A}-1
 \le C\lambda(n+1)e^{C\lambda(n+1)}.
\]
For every lift with a nonzero contribution, its widened-time
bridge has probability at least a fixed positive constant
of staying in the unit tube about its interpolating segment.
Indeed its centered bridge has duration $3\delta\le6$;
Brownian scaling bounds this probability below by the
positive small-ball probability at duration $6$.
The projected tube is contained in $O_n$, so its paths
contribute to the killed kernel on $O_n$.  Summing these
bounds over the endpoint lifts is legitimate: the lifts
give disjoint winding classes in the torus bridge
decomposition.  The cutoffs equal one at every contributing
endpoint.  This yields Equation~\eqref{eq:marked-kernel}.
\end{proof}

The factor $\varepsilon^{1/3}$ records the rarity of the
single-path failure.  The confinement to $O_n$ will ensure
that the interacting moment costs a multiple of $|A|$,
even when the number of marked paths is unrestricted.

\subsection{Colors and the trace-sandwich estimate}

Choose a constant $m_0$ sufficiently large, depending only
on the upper increment constant in Proposition~\ref{prop:increments}
and the fixed $\alpha$-window.  Define
\begin{equation}\label{eq:color-gaps}
 M_n=(m_0+n)/\delta,\qquad \eta_n=l_n e^{m_0+n}.
\end{equation}
We choose $m_0$ so that, for $\delta\le2$,
\begin{equation}\label{eq:gap-choice}
 M_n/2\ge C_{\rm inc}\alpha+1,
 \qquad \beta M_n/2\ge1,
\end{equation}
where $C_{\rm inc}\alpha$ bounds an upper particle-number
increment and $D\ge1$.  The choice of $m_0$ is independent
of $t$.

First truncate the permitted marked bins at $n\le n_*$, so that the
moment in Equation~\eqref{eq:marked-moment-definition} becomes
\[
 \mathfrak M_{n_*}
 =\E_{\cR}\prod_{i=1}^N
 \left[1+\bigl(e^{\lambda J_A(\omega_i)}-1\bigr)
       \ind_E(\omega_i)\ind_{\{n(\omega_i)\le n_*\}}\right].
\]
The enlarged one-particle space is
\[
 \mathfrak h_e=L^2(\Lambda_K)
               \oplus\bigoplus_{n=0}^{n_*}L^2(O_n).
\]
Its first summand is called the original color, and the
others are the marked colors.  In each labelled color
assignment of $N$ particles impose the hard exclusion
only between original-color particles.  The direct sum
of these allowed labelled spaces is denoted by
$\mathcal H_e^{\rm lab}$.  On it let $\mathbb H^{\rm lab}$
be the operator with the original-color hard-core kinetic
form, and with one-particle operator
$-c_1\Delta_{O_n}+M_n$ on a particle of color $n$.
There are no interactions involving a marked color.
Simultaneous permutations of positions and colors commute
with $\mathbb H^{\rm lab}$.  Denote their invariant
subspace by $\mathcal H_e$, and denote the restriction
of $\mathbb H^{\rm lab}$ to it by $\mathbb H$.

On zero-extended original $N$-particle functions apply
in each coordinate the map
\[
 v_e f=(f,(\sqrt{\eta_n}p_n f)_{0\le n\le n_*}),
 \qquad V_{\rm lab}=v_e^{\otimes N}.
\]
Here $V_{\rm lab}$ has domain $L^2(\Omega_N)$, the
original labelled hard-core space, and takes values in
$\mathcal H_e^{\rm lab}$.  It maps form-domain functions
to form-domain functions:
the original hard-core zeros remain, the marked cutoffs
are compactly supported in their $O_n$, and all
multipliers are Lipschitz.  Define first the labelled operator
\[
 \mathcal T_{\rm lab}
 =V_{\rm lab}^*e^{-\delta\mathbb H^{\rm lab}}V_{\rm lab}.
 \]
The map $V_{\rm lab}$ intertwines coordinate permutations
with simultaneous position-color permutations.  Its
restriction $V_e$ to the original symmetric space therefore
takes values in $\mathcal H_e$, and $\mathcal T_{\rm lab}$
commutes with the original symmetrizer $P_{\rm sym}$.
Its symmetric restriction is
\[
 \mathcal T=\mathcal T_{\rm lab}|_{\operatorname{ran}P_{\rm sym}}
           =V_e^*e^{-\delta\mathbb H}V_e.
\]

Expanding the product defining the trajectory moment marks
subsets of labels.  For every marked subset, bin its paths
and drop all collision prohibitions incident to them.
The nonnegative path integral increases.  Apply
Lemma~\ref{lem:marked-kernel} to each marked path.
Its coefficient is reproduced exactly by the color map,
since $\eta_n e^{-\delta M_n}=l_n$.
Each choice of marked labels and their bins appears once
in the labelled color direct sum.  Thus the tested labelled
slab kernel is bounded pointwise by the kernel
$K_{\mathcal T_{\rm lab}}$ of $\mathcal T_{\rm lab}$,
before any restriction to symmetric functions.

The outer kernel $B_{\beta-\delta}^{\rm sym}$ is
nonnegative.  Let $H_N^{\rm lab}$ denote the original
labelled hard-core Hamiltonian.  Integrating the pointwise
bound against this outer kernel gives the symmetric trace
by the following exact identity:
\begin{align*}
 &\int B_{\beta-\delta}^{\rm sym}(Z,Y)
          K_{\mathcal T_{\rm lab}}(Y,Z)\,\dd Y\,\dd Z\\
 &\qquad=\Tr_{\rm lab}\left(
    e^{-(\beta-\delta)H_N^{\rm lab}}
                      P_{\rm sym}\mathcal T_{\rm lab}\right)
 =\Tr_{\rm sym}\left(e^{-(\beta-\delta)H_N}\mathcal T\right).
\end{align*}
Indeed both labelled operators commute with $P_{\rm sym}$;
the product in the middle vanishes on its orthogonal
complement and restricts to the final product on its range.
The sole permutation normalization is the average
$P_{\rm sym}=N!^{-1}\sum_\pi U_\pi$ already present
in the outer kernel.  No factorial is inserted in the
labelled color expansion or in this trace identity.
Consequently the truncated moment satisfies
\begin{equation}\label{eq:slab-holder}
 \mathfrak M_{n_*}
 \le\frac{\Tr_{\rm sym}
          (e^{-(\beta-\delta)H_N}\mathcal T)}{Z_N}
 \le\left(\frac{\Tr_{\rm sym}\mathcal T^{\beta/\delta}}
                   {Z_N}\right)^{\delta/\beta}.
\end{equation}
The second inequality is Schatten H\"older with exponents
$\beta/(\beta-\delta)$ and $\beta/\delta$.
All operators have finite positive-time traces at the
fixed finite volume and particle number.

To control the large power in Equation~\eqref{eq:slab-holder},
we must account for the norm of the color map.  Define
\begin{equation}\label{eq:color-normalization}
 a(x)=\frac12\log\left(1+\sum_{n=0}^{n_*}\eta_np_n(x)^2\right),
 \qquad A_0=\sum_{i=1}^N a(x_i).
\end{equation}
The same position multiplier acts in every color.
The map $U_e=e^{-A_0}V_e$ is isometric, because the
squared pointwise norm of its one-coordinate color vector
is one.  Therefore
\[
 \mathcal T=U_e^*X_0U_e,
 \qquad X_0=e^{A_0}e^{-\delta\mathbb H}e^{A_0}.
\]
Eigenvalue min--max for positive compact compressions gives
$\Tr\mathcal T^q\le\Tr X_0^q$ for every $q>0$.
The next lemma converts this sandwich into a heat trace
with a bounded attractive field.

\begin{lemma}[Trace bound for a Lipschitz sandwich]\label{lem:sandwich}
For the preceding finite-color space and operators, set
\[
 W_0=2a/\delta+c_1|\nabla a|^2,\qquad
 \widetilde{\mathbb H}=\mathbb H-\sum_{i=1}^N W_0(x_i).
\]
For every real $q>0$,
\begin{equation}\label{eq:all-q-sandwich}
 \Tr_{\mathcal H_e}X_0^q
 \le\Tr_{\mathcal H_e}e^{-q\delta\widetilde{\mathbb H}}.
\end{equation}
\end{lemma}

\begin{proof}
This is the sandwich estimate of \citet[Lemma ``Trace bound
for a Lipschitz sandwich'']{B}, with the hard-core
Dirichlet conditions retained.  We give the form argument,
including the step that permits every positive power.

There are finitely many color assignments.  In each,
Dirichlet domain monotonicity bounds the heat trace by a
free periodic or Dirichlet product heat trace.  Hence
$\mathbb H$ has compact resolvent and finite heat traces
at positive times.  Subtracting the bounded potential
$\sum_i W_0(x_i)$ preserves these properties.
The operator $X_0$ is positive, compact, and injective.

For $0\le s\le\delta$ define
\[
 U(s)=e^{u(s)A_0}e^{-s\mathbb H}e^{A_0},
 \qquad u(s)=-1+2s/\delta.
\]
Thus $U(0)=I$ and $U(\delta)=X_0$.
Bounded Lipschitz multiplication preserves each form
domain, including its original-color collision boundary
and marked-color spatial boundary.  This follows by
applying the product rule to compactly supported
approximants and taking their form-norm limits.
For an initial vector $f_0$ and $s>0$, spectral calculus
puts $e^{-s\mathbb H}e^{A_0}f_0$ in the operator domain.
The vector $f=U(s)f_0$ is norm differentiable and in the
form domain.  If $\mathfrak q$ is the form of
$\mathbb H$, then
\begin{equation}\label{eq:sandwich-derivative}
 \frac12\frac{\dd}{\dd s}\|f\|^2
 =\frac2\delta\langle f,A_0f\rangle
 -\operatorname{Re}\mathfrak q(e^{u(s)A_0}f,e^{-u(s)A_0}f).
\end{equation}
This identity only pairs a form-domain vector with the
operator-domain semigroup vector; it does not require
Lipschitz multiplication to preserve the operator domain.

In any color assignment, the kinetic coefficient $c_i$
of coordinate $i$ is $1$ or $c_1$.  The product rule
gives, almost everywhere,
\[
 \operatorname{Re}\left[
 \overline{\nabla_i(e^{uA_0}f)}\cdot
                  \nabla_i(e^{-uA_0}f)\right]
 =|\nabla_i f|^2-u^2|\nabla a(x_i)|^2|f|^2.
\]
The cross terms have zero real part and the gap
potentials commute with the multipliers.  Since
$|u|\le1$ and $c_i\le c_1$, Equation~\eqref{eq:sandwich-derivative}
implies
\begin{equation}\label{eq:sandwich-form}
 \frac12\frac{\dd}{\dd s}\|U(s)f_0\|^2
 \le-\widetilde{\mathfrak q}[U(s)f_0],
\end{equation}
where $\widetilde{\mathfrak q}$ is the form of
$\widetilde{\mathbb H}$.

Let $\widetilde E_1\le\widetilde E_2\le\cdots$
be the eigenvalues of $\widetilde{\mathbb H}$ with
multiplicity.  Apply the same form calculation to
$U(s)^{\otimes k}$, using the sums of $A_0$ and
$\mathbb H$ over $k$ copies of $\mathcal H_e$.
Restrict these copies to their alternating subspace.
In an eigenbasis of $\widetilde{\mathbb H}$, the
copy-sum Hamiltonian on that subspace has lower bound
$\sum_{i=1}^k\widetilde E_i$.  The copies here are
copies of the entire enlarged symmetric $N$-particle
space; the particle symmetry within each copy is
unchanged.  Equation~\eqref{eq:sandwich-form} consequently
gives
\[
 \|\mathop{\bigwedge}\nolimits^k U(\delta)\|
 \le\exp\left(-\delta\sum_{i=1}^k\widetilde E_i\right).
\]
To justify integration at zero, integrate first from
$s=\epsilon>0$ and then let $\epsilon\downarrow0$,
using strong continuity and $U(0)=I$.  The tensor
form calculation is first valid on finite tensor
sums and extends by the closed forms and density.

Write the decreasing eigenvalues of $X_0$ as $x_i>0$.
The norm of its $k$th exterior power is the product of
its largest $k$ eigenvalues.  Hence
\[
 \sum_{i=1}^k\log x_i
 \le-\delta\sum_{i=1}^k\widetilde E_i
 \qquad(k\ge1).
\]
For each real $h$, taking the supremum over $k\ge0$
of the shifted partial sums proves
\[
 \sum_i(\log x_i-h)_+
 \le\sum_i(-\delta\widetilde E_i-h)_+.
\]
Both sums have finitely many positive terms for fixed
$h$, by compactness.  For every $q>0$,
\[
 e^{qu}=\int_{\mathbb R}q^2e^{qh}(u-h)_+\,\dd h.
\]
Integrate the preceding positive-part inequality and
apply Tonelli's theorem.  It gives
$\sum_i x_i^q\le\sum_i e^{-q\delta\widetilde E_i}$,
which is Equation~\eqref{eq:all-q-sandwich}.
\end{proof}

\subsection{The interacting trajectory moment}

The remaining task is to compare the enlarged heat trace
with $Z_N$.  The original color will be estimated by the
field-density bound; the positive gaps will pay for the
other colors.  Spatial confinement in
Equation~\eqref{eq:bin-volume} keeps both costs
proportional to the number of tested sites.

\begin{proposition}[Joint trajectory moment]\label{prop:trajectory-moment}
Use any of the three events $E$ and the coefficients
$l_n,\eta_n,M_n$ from Equations~\eqref{eq:marked-kernel}
and \eqref{eq:color-gaps}.  For a sufficiently large
fixed constant $C_0$, assume
\begin{equation}\label{eq:color-smallness}
 \frac{C_0}{\delta}\sum_{n\ge0}\eta_n\le1.
\end{equation}
Then for every nonempty set $A$ of sites,
\begin{align}
 &\log\E_{\cR}\exp\left(
       \lambda\sum_{i=1}^N J_A(\omega_i)\ind_E(\omega_i)\right)
 \notag\\
 &\quad\le
 CD|A|\sum_{n\ge0}(n+1)^3\eta_n
 +C\delta T|A|(V^{-1}+T^{3/2})
          \sum_{n\ge0}(n+1)^3e^{-\beta M_n/2}.
 \label{eq:trajectory-moment}
\end{align}
All constants are independent of $n_*,K,D,T$ and
$0<t\le1$, with the stated $t$-dependence retained in
$\varepsilon$ for irregularity.
\end{proposition}

\begin{proof}
For the truncated normalization in
Equation~\eqref{eq:color-normalization}, write
$S(x)=\sum_{n\le n_*}\eta_np_n(x)^2$.
Cauchy--Schwarz gives
\[
 |\nabla a|^2
 \le\frac{S}{(1+S)^2}
               \sum_{n\le n_*}\eta_n|\nabla p_n|^2
 \le C\sum_{n\le n_*}\eta_n\ind_{O_n}.
\]
Also $2a\le S$.  Since $\delta\le2$,
\begin{equation}\label{eq:color-field-bound}
 0\le W_0\le\frac C\delta
                  \sum_{n\le n_*}\eta_n\ind_{O_n}\le1,
\end{equation}
where the last inequality follows from
Equation~\eqref{eq:color-smallness}.

For every $0\le J\le N$, differentiation with respect
to a bounded field strength gives
\[
 \frac{\dd}{\dd u}\log Z_J^{(-uW_0)}
 =\beta\E_{J,-uW_0}\sum_{i=1}^J W_0(x_i)
 \qquad(0\le u\le1).
\]
At finite volume this identity follows from the
Duhamel formula and trace cyclicity.  The field is
bounded and the positive-time traces are finite.
Apply Corollary~\ref{cor:field-density} to the cell
covers in Equation~\eqref{eq:bin-volume}, and then
integrate in $u$.  It follows that
\begin{equation}\label{eq:original-color-field}
 \log\frac{Z_J^{(-W_0)}}{Z_J}
 \le\frac{\beta C}{\delta}D|A|
                     \sum_{n\le n_*}(n+1)^3\eta_n.
\end{equation}
The vacuum case has zero left side.  In addition,
Proposition~\ref{prop:increments} gives
\begin{equation}\label{eq:color-removal-cost}
 \frac{Z_J}{Z_N}\le e^{\beta C_{\rm inc}\alpha(N-J)}.
\end{equation}

Evaluate Lemma~\ref{lem:sandwich} at $q=\beta/\delta$.
The enlarged symmetric space decomposes by the number
$J$ of original-color particles and the occupation
numbers $k_n$ of each marked color, with
$J+\sum_n k_n=N$.  Each sector is one normalized
tensor product of the symmetric species spaces.
There is no multinomial multiplicity: symmetrizing
the mutually orthogonal assignments of colors gives
an isometry from this tensor product onto the sector.
Its original factor has trace $Z_J^{(-W_0)}$;
the marked factors are ideal Bose canonical traces
on $O_n$.

Use Equations~\eqref{eq:original-color-field} and
\eqref{eq:color-removal-cost}.  Absorbing the latter
factor into the marked energies subtracts
$C_{\rm inc}\alpha$ per marked particle.  With
$W_0\le1$ and Equation~\eqref{eq:gap-choice}, each
such one-particle operator is bounded below in form
by $-c_1\Delta_{O_n}+M_n/2$.  The variational
eigenvalue principle therefore bounds every fixed
marked canonical trace by that operator's trace.
We may then drop the restriction $\sum_nk_n\le N$
and sum the independent marked occupations.  The
original count is determined by those occupations
before this restriction is dropped, so no extra
sum over $J$ is introduced.

The logarithm of the ideal grand trace for color
$n$ is
\[
 \sum_{a\ge1}\frac{e^{-a\beta M_n/2}}a
                  \Tr e^{a\beta c_1\Delta_{O_n}}
 \le C|O_n|(V^{-1}+T^{3/2})e^{-\beta M_n/2}.
\]
Here the periodic Gaussian sum and killing give
$\Tr e^{s c_1\Delta_{O_n}}
 \le C|O_n|(V^{-1}+s^{-3/2})$,
and $\beta M_n/2\ge1$ controls the sum in $a$.
Combining this with the two preceding comparisons
proves
\begin{align*}
 \log\frac{\Tr e^{-\beta\widetilde{\mathbb H}}}{Z_N}
 &\le\frac{\beta C}{\delta}D|A|
                      \sum_{n\le n_*}(n+1)^3\eta_n\\
 &\quad+C|A|(V^{-1}+T^{3/2})
                      \sum_{n\le n_*}(n+1)^3e^{-\beta M_n/2}.
\end{align*}
Multiply by $\delta/\beta=\delta T$ and use
Equation~\eqref{eq:slab-holder}, the compression
inequality, and Lemma~\ref{lem:sandwich}.
This proves Equation~\eqref{eq:trajectory-moment}
for the truncated product.

For each path the truncated marking weight increases
to the full marking weight as $n_*\to\infty$.
The finite product over labels also increases.
Monotone convergence removes the truncation, while
the displayed bounds on the right are uniform in
$n_*$ and summable.  This proves the proposition.
\end{proof}

\subsection{Completion of the certificate estimate}

We first record the sizes of the two costs in
Equation~\eqref{eq:trajectory-moment}.  The Gaussian
bin tail in Equation~\eqref{eq:marked-kernel} implies
\begin{equation}\label{eq:eta-sum}
 \sum_{n\ge0}(n+1)^3\eta_n
 \le C_\lambda\lambda\varepsilon^{1/3},
\end{equation}
where $C_\lambda$ is bounded on bounded
$\lambda$-intervals, uniformly for $0<t\le1$.
Indeed the summand is bounded by a constant times
$\lambda\varepsilon^{1/3}(n+1)^4
 e^{C\lambda(n+1)+n-cn^2}$, including the fixed
factor $e^{m_0}$.  The same estimate without
$(n+1)^3$ controls the smallness condition.
Since $\beta\ge D$ and $\delta\le2$, the second
term of Equation~\eqref{eq:trajectory-moment} is
$o(1)|A|$ as $D\to\infty$, uniformly in
$0<t\le1$ and $K\ge100$.

\begin{proof}[Proof of Proposition~\ref{prop:certificates}]
It suffices to prove arbitrarily small exponential
bases for simultaneous failures of each of the six
individual tests.  At both endpoints the law is
$\mu$, by Proposition~\ref{prop:purification}.
Thus the low-count test at $Y$ and the close-count
tests at $Y,Z$ have this property by
Propositions~\ref{prop:low-count} and \ref{prop:close-count}.

For the number of visitors choose $E=E_{\rm all}$
and $\lambda=c't$, where $c'>0$ is sufficiently
small and fixed.  Equation~\eqref{eq:eta-sum}
makes Equation~\eqref{eq:color-smallness} hold
uniformly for $0<t\le1$, and the moment's first
coefficient is at most $C\lambda D$ with $C$
independent of $t$.  If every site of $A$ has more
than $H_1D$ visitors, then
$\sum_i J_A(\omega_i)>H_1D|A|$.  Exponential
Markov inequality yields
\[
 \cR(\text{this event})
 \le\exp\{[-\lambda(H_1-C)D+o(1)]|A|\}.
\]
Choose $H_1>C+1$ now.  For every subsequently
fixed $t>0$, the exponential base tends to zero
as $D$ increases.  This choice of $H_1$ is
unaffected by reducing $t$.

For large-diameter visitors take $\lambda=1$ and
$E=E_{\rm diam}$.  Now
$\varepsilon^{1/3}\le Ce^{-c/t}$.
Choose $t$ sufficiently small that both
Equation~\eqref{eq:color-smallness} and the bound
$CD\sum_n(n+1)^3\eta_n\le\theta D/2$ hold.
If every site of $A$ has at least $\theta D$
such visitors, their total number of visits to
the tested balls is at least $\theta D|A|$.
Equation~\eqref{eq:trajectory-moment} therefore
bounds this probability by
$\exp[-(\theta D/2-o(1))|A|]$.

Finally fix a large $\lambda>0$ and take
$E=E_{\rm irr}$.  With $t$ now fixed,
Equation~\eqref{eq:eta-sum} gives
$\sum_n(n+1)^3\eta_n=O_{\lambda,t}(D^{-4})$.
The smallness condition holds for large $D$,
and the whole coefficient per site in the
moment tends to zero.  If each tested site has
even one irregular visitor, the sum of its
irregular visits is at least $|A|$, counting
a path once for every tested ball it visits.
Its probability is consequently at most
$\exp[-(\lambda-o(1))|A|]$.
Choosing $\lambda$ large and then $D$ large
makes this base as small as prescribed.

For clarity, no independence between tests or
between cells is needed to combine them.
Given a failure at every site, assign one of
the six failure types to each site.  There
are at most $6^{|A|}$ assignments.  In each,
some type occurs on a prescribed subset of
size at least $|A|/6$.  If every simultaneous
single-type failure has probability at most
$(p')^{|A'|}$, with $0<p'<1$, the union is
bounded by $(6(p')^{1/6})^{|A|}$.
Take $p'$ so small that $6(p')^{1/6}\le p_1$.
The preceding choices and a final increase
of $D_0$ ensure all six estimates.  The empty
set satisfies Equation~\eqref{eq:joint-certificates}
trivially.  All arguments survive a further
reduction of the fixed $t$, followed by the
corresponding increase of $D_0$.
\end{proof}

We have obtained an unconditional joint bound for
failure of the full-path certificates, with no loss
proportional to the total volume.  The next two
sections use these certificates to find movable
paths and then expose their environment while
preserving an exact conditional bridge law.
